%
% Chapter 5 — Data Model
%
\chapter{Data Model} \label{cap:datamodel}

\section{Entities} \label{sec:dm_entities}

This section describes the core entities of the Diário LX data model, their
attributes, and the relationships between them. The complete entity-relationship
diagram is presented in Appendix~\ref{app:db_model}.

\subsection{User} \label{sec:dm_user}

Represents an authenticated system user. Each user has a \texttt{user\_role}
(\textit{ADMIN}, \textit{EDITOR}, or \textit{CONTRIBUTOR}), credentials
(\texttt{username}, \texttt{email}, \texttt{password\_hash}), profile information
(\texttt{first\_name}, \texttt{last\_name}, \texttt{bio}), and an optional
\texttt{avatar\_media\_id} reference to a media item. An \texttt{active\_account}
flag is used to deactivate access without deleting the record.

\subsection{Session} \label{sec:dm_session}

Stores active refresh tokens issued upon login. Each session is associated with a
user and holds a \texttt{refresh\_token}, \texttt{created\_at}, and
\texttt{expires\_at}. Sessions are deleted on logout and expire automatically based
on \texttt{expires\_at}.

\subsection{Invite} \label{sec:dm_invite}

Represents an invitation token generated by an Administrator to allow new user
registration. Each invite carries a \texttt{role\_assigned}, an expiration
timestamp, and a \texttt{used} flag to prevent reuse.

\subsection{Content} \label{sec:dm_content}

The central entity of the system. Each content item has a \texttt{content\_type}
(\textit{ARTICLE}, \textit{VIDEO}, \textit{EPISODE}, or \textit{PODCAST}), a
\texttt{title}, a \texttt{headline}, a \texttt{slug}, and a reference to its
\texttt{category}. The body of a content item is composed of an ordered sequence
of \texttt{content\_blocks}.

The \texttt{published\_at} timestamp controls content visibility. When set to a
past value, the content is immediately publicly accessible. When set to a future
value, the content is scheduled for publication --- it has been approved but will
only become visible once that timestamp is reached. A \texttt{NULL} value indicates
the content has not yet been approved for publication.

Content items may have one or more authors, recorded in the \texttt{content\_authors}
association table with a \texttt{role} of \textit{primary} or \textit{secondary}.
A partial unique index ensures that each content item has at most one primary author.

\subsection{Content Block} \label{sec:dm_block}

Represents a single unit of content within a content item. Each block has a
\texttt{type} (\textit{text}, \textit{quote}, or \textit{image}), a
\texttt{position} defining its order within the content item, and either a
\texttt{content} text field (for text and quote blocks) or a \texttt{media\_id}
reference (for image blocks). A database constraint enforces that text and quote
blocks must have \texttt{content} set, and image blocks must have \texttt{media\_id}
set.

\subsection{Category} \label{sec:dm_category}

Editorial categories support a two-level hierarchy via a \texttt{parent\_id}
self-reference on the \texttt{categories} table. Top-level categories have
\texttt{parent\_id = NULL}; subcategories reference their parent. Each category
has a \texttt{name}, a unique \texttt{slug}, a unique \texttt{color}, and an
optional \texttt{description}. Categories may be soft-deleted via an
\texttt{archived\_at} timestamp.

\subsection{Tag} \label{sec:dm_tag}

Tags provide cross-cutting thematic classification. A content item may have multiple
tags, recorded in the \texttt{content\_tags} association table. Each tag association
carries a \texttt{role} of \textit{primary} or \textit{secondary}; a partial unique
index ensures that each content item has at most one primary tag.

\subsection{Media} \label{sec:dm_media}

Stores metadata for uploaded media files (images, videos, and audio). Binary files
are not stored in the database --- only their \texttt{bucket}, \texttt{object\_key},
\texttt{mime\_type}, \texttt{size\_bytes}, and a \texttt{status} field
(\textit{pending} or \textit{ready}) are persisted. The \textit{pending} state
indicates that the frontend has initiated the upload but not yet confirmed it; the
backend transitions the status to \textit{ready} upon confirmation.

Media items may have associated credits via the \texttt{media\_credits} table,
which records the contributing user and their \texttt{credit\_role}
(e.g.\ \textit{PHOTOGRAPHER}, \textit{VIDEOGRAPHER}, \textit{AUDIOGRAPHER}).

\subsection{Featured Content} \label{sec:dm_featured}

The \texttt{featured\_contents} table stores curated content selections for the
homepage. Each entry has a \texttt{key} identifying the slot (e.g.\
\texttt{homepage\_hero}, \texttt{lisboa\_featured}), a mandatory \texttt{content\_id}
reference to the selected content item, and a \texttt{position} for ordering within
the same slot group.

\section{Database Views} \label{sec:dm_views}

Several PostgreSQL views are defined to encapsulate complex query logic and provide
pre-joined representations of the data for use by the repository layer:

\begin{description}
      \item[\texttt{v\_contents}] Full content view, including category, featured media,
            authors ordered by role, tags, and content blocks with their associated media.
            Used for single content item retrieval.

      \item[\texttt{v\_contents\_summary}] Lightweight content listing view, including
            only the fields required for listing pages: title, slug, type, category,
            featured image URL, primary tag, and author names. Used for paginated content
            listings.

      \item[\texttt{v\_medias}] Media view that computes the full \texttt{url} and
            \texttt{thumbnail\_url} from \texttt{bucket} and \texttt{object\_key}, and
            aggregates credits as a JSON array.

      \item[\texttt{v\_categories}] Category view that joins each category with its
            parent name and includes a \texttt{quantity} count of associated content items.

      \item[\texttt{v\_tags}] Tag view that includes a \texttt{quantity} count of
            associated content items.
\end{description}

\section{Indexes} \label{sec:dm_indexes}

The following indexes were defined to support the most frequent query patterns:

\begin{description}
      \item[\texttt{idx\_content\_category\_id}] Supports filtering content by category.

      \item[\texttt{idx\_content\_blocks\_content\_position}] Composite index on
            \texttt{(content\_id, position)}, supporting ordered block retrieval per
            content item.

      \item[\texttt{idx\_content\_authors\_content\_id} and \texttt{idx\_content\_authors\_author\_id}]
            Support join queries on the \texttt{content\_authors} association table.

      \item[\texttt{uq\_content\_primary\_author}] Partial unique index ensuring that
            each content item has at most one author with \texttt{role = 'primary'}.

      \item[\texttt{uq\_content\_primary\_tag}] Partial unique index ensuring that
            each content item has at most one tag with \texttt{role = 'primary'}.
\end{description}

\section{Key Design Decisions} \label{sec:dm_decisions}

\begin{itemize}
      \item \textbf{Block-based content model}: Instead of storing article body as a
            single rich-text blob, content is decomposed into an ordered sequence of
            typed blocks (\textit{text}, \textit{quote}, \textit{image}). This enables
            structured rendering, easier reordering, and per-block media references.

      \item \textbf{Self-referencing category hierarchy}: A single \texttt{categories}
            table with a \texttt{parent\_id} foreign key supports the two-level category
            hierarchy without requiring a separate \texttt{subcategories} table.

      \item \textbf{Media metadata separation}: Binary files are stored in SeaweedFS;
            only metadata is persisted in the database. The \textit{pending}/\textit{ready}
            status tracks the asynchronous upload lifecycle, ensuring that incomplete
            uploads do not surface in the application before the file is available in
            storage.

      \item \textbf{Timestamps as epoch integers}: All timestamps are stored as
            \texttt{BIGINT} Unix epoch values rather than \texttt{TIMESTAMP} columns.
            This avoids timezone conversion issues, as epoch values are inherently UTC,
            and simplifies interoperability with the \texttt{kotlinx.datetime.Instant}
            type used throughout the backend.

      \item \textbf{Serial primary keys}: All tables use \texttt{SERIAL}
            (auto-incrementing integer) primary keys rather than UUIDs. This simplifies
            joins and keeps index sizes small, which is well-suited to the expected data
            volume of the platform.
\end{itemize}